\section{Análisis comparativo entre LangChain y LangGraph}

\begin{table}[H]
\centering
\begin{tabular}{lll}
\toprule
\textbf{Dimensión} & \textbf{LangChain} & \textbf{LangGraph} \\
\midrule
Gestión del historial & Texto añadido a un solo mensaje & Lista de mensajes de varios turnos \\
Estructura de mensajes & Un solo mensaje human (con scratchpad) & Secuencia de mensajes Human / AI / Tool \\
Llamadas a herramientas & En serie (una a la vez) & Con soporte para paralelismo (lista de tool\_calls) \\
Arquitectura & Plantilla de prompt + Executor & Graph (nodos + bordes condicionales) \\
Visualización & Requiere seguimiento manual & Se puede exportar directamente el diagrama del Graph \\
\bottomrule
\end{tabular}
\caption{Comparación de la implementación de ReAct en LangChain y LangGraph}
\end{table}

\subsection{Resumen de la sección}
LangGraph supera de manera integral a la implementación anterior de LangChain en claridad de arquitectura, gestión de mensajes y capacidad de paralelismo. El enfoque de LangChain resulta más adecuado para entender el concepto original de ReAct, mientras que LangGraph es más apropiado para entornos de producción.

